\documentclass[aspectratio=169]{beamer}

% ── Theme & colours ──────────────────────────────────────────────────────────
\usetheme{Madrid}
\usecolortheme{whale}
\usefonttheme{professionalfonts}
\usepackage{array}
\definecolor{primary}{RGB}{0,74,128}
\definecolor{accent}{RGB}{0,153,204}
\definecolor{lightgray}{RGB}{245,245,245}

\setbeamercolor{palette primary}{bg=primary,fg=white}
\setbeamercolor{palette secondary}{bg=accent,fg=white}
\setbeamercolor{palette tertiary}{bg=primary!80,fg=white}
\setbeamercolor{palette quaternary}{bg=primary,fg=white}
\setbeamercolor{structure}{fg=primary}
\setbeamercolor{section in toc}{fg=primary}
\setbeamercolor{block title}{bg=primary,fg=white}
\setbeamercolor{block body}{bg=lightgray,fg=black}
\setbeamercolor{alerted text}{fg=accent}

% ── Packages ─────────────────────────────────────────────────────────────────
\usepackage[utf8]{inputenc}
\usepackage[T1]{fontenc}
\usepackage[french]{babel}
\usepackage{lmodern}
\usepackage{microtype}
\usepackage{graphicx}
\usepackage{booktabs}
\usepackage{tikz}
\usetikzlibrary{positioning,shapes.geometric,arrows.meta,fit}
\usepackage{fontawesome5}
\usepackage{xcolor}
\usepackage{enumitem}

% ── Remove navigation symbols ────────────────────────────────────────────────
\setbeamertemplate{navigation symbols}{}

% ── Footer ───────────────────────────────────────────────────────────────────
\setbeamertemplate{footline}{%
  \leavevmode%
  \hbox{%
    \begin{beamercolorbox}[wd=.33\paperwidth,ht=2.25ex,dp=1ex,center]{palette primary}%
      \usebeamerfont{author in head/foot}\insertshortauthor
    \end{beamercolorbox}%
    \begin{beamercolorbox}[wd=.34\paperwidth,ht=2.25ex,dp=1ex,center]{palette secondary}%
      \usebeamerfont{title in head/foot}Alexandre — Microservices E-commerce
    \end{beamercolorbox}%
    \begin{beamercolorbox}[wd=.33\paperwidth,ht=2.25ex,dp=1ex,center]{palette primary}%
      \usebeamerfont{date in head/foot}%
      \insertframenumber{} / 14
    \end{beamercolorbox}%
  }%
  \vskip0pt%
}

% ── Bullet style ─────────────────────────────────────────────────────────────
\setbeamertemplate{itemize item}{\color{accent}\faAngleRight}
\setbeamertemplate{itemize subitem}{\color{primary}\faAngleRight}

% ── Meta ─────────────────────────────────────────────────────────────────────
\title[Alexandre — PFA]%
      {\textbf{Alexandre}\\[4pt]
       \large Une application e-commerce en microservices}

\author[Nom de l'étudiant]{%
  \textbf{[Nom de l'étudiant]}\\[4pt]
  \small CI2 — Génie Informatique
}

\institute[Jamji Express]{%
  Projet de Fin d'Année\\[2pt]
  Entreprise d'accueil : \textbf{Jamji Express}
}

\date{Année universitaire [Année]}

% ═══════════════════════════════════════════════════════════════════════════════
\begin{document}
% ═══════════════════════════════════════════════════════════════════════════════

% ── Slide 1 — Titre ──────────────────────────────────────────────────────────
\begin{frame}[plain]
  \titlepage
\end{frame}

% ── Slide 2 — Organisme d'accueil ────────────────────────────────────────────
\begin{frame}{Présentation de l'organisme d'accueil}

  \begin{block}{\faBuilding\quad Jamji Express}
    Entreprise spécialisée dans l'\textbf{import et la logistique},
    assurant la gestion et l'acheminement de marchandises.
  \end{block}

  \vspace{1em}

  \begin{columns}[T]
    \begin{column}{0.48\textwidth}
      \textbf{\color{primary}Contexte du stage}
      \begin{itemize}
        \item Stage en milieu professionnel réel
        \item Liberté de choix du sujet accordée à l'étudiant
        \item Orientation vers les \alert{architectures logicielles modernes}
      \end{itemize}
    \end{column}
    \begin{column}{0.48\textwidth}
      \textbf{\color{primary}Démarche adoptée}
      \begin{itemize}
        \item Exploration des \alert{microservices} en pratique
        \item Conception d'un projet complet de bout en bout
        \item Mise en œuvre des bonnes pratiques industrielles
      \end{itemize}
    \end{column}
  \end{columns}

\end{frame}

% ── Slide 3 — Contexte et Motivation ─────────────────────────────────────────
\begin{frame}{Contexte et Motivation}

  \begin{block}{\faExclamationTriangle\quad Les exigences du logiciel moderne}
    Les applications d'aujourd'hui exigent \textbf{scalabilité},
    \textbf{maintenabilité} et \textbf{déploiement continu}.
  \end{block}

  \vspace{0.8em}

  \begin{columns}[T]
    \begin{column}{0.48\textwidth}
      \textbf{\color{primary}Une solution éprouvée}
      \begin{itemize}
        \item L'architecture microservices s'est imposée comme réponse industrielle
        \item Adoptée par \alert{Amazon}, \alert{Netflix}, \alert{Uber}
        \item Découplage, autonomie, résilience
      \end{itemize}
    \end{column}
    \begin{column}{0.48\textwidth}
      \textbf{\color{primary}Objectif du projet}
      \begin{itemize}
        \item Explorer cette architecture \alert{en pratique}
        \item Terrain d'expérimentation technique
        \item Application e-commerce simplifiée, non destinée à la production
      \end{itemize}
    \end{column}
  \end{columns}

  \vspace{0.8em}
  \begin{alertblock}{}
    \centering
    \small La motivation principale est \textbf{l'apprentissage par la pratique}
    des concepts fondamentaux des systèmes distribués.
  \end{alertblock}

\end{frame}

% ── Slide 4 — Problématique ───────────────────────────────────────────────────
\begin{frame}{Problématique}

  \begin{center}
    \begin{tikzpicture}
      \node[draw=primary, fill=primary!10, rounded corners=6pt,
            text width=0.85\textwidth, align=center,
            inner sep=12pt, font=\large\itshape] (q) {
        « Comment concevoir et implémenter une application web selon une
        architecture microservices, en garantissant une
        \textbf{séparation claire des responsabilités},
        une \textbf{communication efficace entre les services},
        et une \textbf{facilité de déploiement} ? »
      };
    \end{tikzpicture}
  \end{center}

  \vspace{0.6em}

  \begin{columns}[T]
    \begin{column}{0.32\textwidth}
      \centering
      {\color{primary}\faLayerGroup}\\[4pt]
      \textbf{Séparation}\\
      \small des responsabilités par service
    \end{column}
    \begin{column}{0.32\textwidth}
      \centering
      {\color{primary}\faNetworkWired}\\[4pt]
      \textbf{Communication}\\
      \small synchrone et asynchrone
    \end{column}
    \begin{column}{0.32\textwidth}
      \centering
      {\color{primary}\faRocket}\\[4pt]
      \textbf{Déploiement}\\
      \small continu et reproductible
    \end{column}
  \end{columns}

  \vspace{0.8em}
  \begin{block}{}
    \centering
    \small Ce projet constitue une \textbf{réponse concrète} à cette question,
    en mettant en pratique les concepts des architectures distribuées modernes.
  \end{block}

\end{frame}

% ── Slide 5 — Présentation du Projet ─────────────────────────────────────────
\begin{frame}{Présentation du Projet — Alexandre}

  \begin{block}{\faShoppingCart\quad Qu'est-ce qu'Alexandre ?}
    Plateforme e-commerce basée sur une architecture microservices, développée avec
    \textbf{Spring Boot} (Java), conçue comme \textbf{projet d'apprentissage} et de
    démonstration des patterns d'architecture distribuée moderne.
  \end{block}

  \vspace{0.6em}

  \begin{columns}[T]
    \begin{column}{0.48\textwidth}
      \textbf{\color{primary}Patterns illustrés}
      \begin{itemize}
        \item \faSearch\; \textbf{Service Discovery} (Eureka)
        \item \faCog\; \textbf{Configuration centralisée}
        \item \faBell\; \textbf{Notifications event-driven} (Kafka)
        \item \faChartLine\; \textbf{Observabilité complète}
              \begin{itemize}
                \item Logs (Loki)
                \item Métriques (Prometheus)
                \item Tracing (Tempo / OpenTelemetry)
              \end{itemize}
      \end{itemize}
    \end{column}
    \begin{column}{0.48\textwidth}
      \textbf{\color{primary}Perspectives d'évolution}
      \begin{itemize}
        \item Communication \alert{asynchrone} généralisée
        \item \alert{Réplicas horizontaux} et autoscaling
        \item \alert{Circuit breakers} (Resilience4j)
        \item Mise en \alert{cache} (Redis)
      \end{itemize}

      \vspace{0.6em}
      \begin{alertblock}{\small Bilan}
        \small Projet \textbf{structuré et extensible}, posant des
        fondations solides pour une mise à l'échelle ciblée.
      \end{alertblock}
    \end{column}
  \end{columns}

\end{frame}

% ── Slide 6 — Pourquoi les Microservices ? ───────────────────────────────────
\begin{frame}{Pourquoi les Microservices ?}

  \begin{center}
    \renewcommand{\arraystretch}{1.35}
    \begin{tabular}{>{\bfseries\color{primary}}l
                    >{\centering\arraybackslash}p{0.28\textwidth}
                    >{\centering\arraybackslash}p{0.28\textwidth}}
      \toprule
      \textbf{\color{black}Critère}
        & \textbf{Monolithe}
        & \textbf{Microservices} \\
      \midrule
      Déploiement
        & Un seul artefact
        & Indépendant par service \\
      Scalabilité
        & Globale uniquement
        & Ciblée par service \\
      Stack technique
        & Unique et partagée
        & Polyglotte possible \\
      Tolérance aux pannes
        & Panne globale
        & Isolation des défaillances \\
      Maintenabilité
        & Couplage fort
        & Découplage, responsabilité unique \\
      Complexité initiale
        & \textcolor{accent}{\textbf{Faible}}
        & \textcolor{red!70!black}{\textbf{Élevée dès le départ}} \\
      \bottomrule
    \end{tabular}
  \end{center}

  \vspace{0.4em}

  \begin{alertblock}{}
    \small Les microservices introduisent une \textbf{complexité opérationnelle réelle} —
    leur adoption est justifiée à grande échelle ou dans un
    \textbf{objectif d'apprentissage explicite}.
  \end{alertblock}

\end{frame}

% ── Slide 7 — Architecture globale ───────────────────────────────────────────
\begin{frame}{Architecture Globale d'Alexandre}

  \begin{center}
    \includegraphics[width=\textwidth,height=0.30\textheight,keepaspectratio]{architecture.jpeg}
  \end{center}

  \vspace{0.2em}

  \begin{columns}[T]
    \begin{column}{0.30\textwidth}
      \textbf{\color{primary}\small Couche d'entrée}
      \begin{itemize}\small
        \item \textbf{Nginx} — reverse proxy
        \item \textbf{API Gateway}
        \item \textbf{Eureka} — service discovery
        \item \textbf{Keycloak} — auth. JWT
      \end{itemize}
    \end{column}
    \begin{column}{0.38\textwidth}
      \textbf{\color{primary}\small Services métier}
      \begin{itemize}\small
        \item User, Order, Inventory
        \item Payment, Delivery
        \item \textbf{Notification} (event-driven / Kafka)
        \item \textbf{Config Server}
        \item[$\hookrightarrow$] \textit{\scriptsize Chaque service : base de données dédiée}
      \end{itemize}
    \end{column}
    \begin{column}{0.28\textwidth}
      \textbf{\color{primary}\small Observabilité}
      \begin{itemize}\small
        \item \textbf{Prometheus} — métriques
        \item \textbf{Loki} — logs
        \item \textbf{Tempo} — tracing distribué
        \item \textbf{Grafana} — visualisation
      \end{itemize}
    \end{column}
  \end{columns}

  \begin{block}{}
    \small \textbf{Base de données par service} — principe fondamental garantissant
    un découplage réel entre les services.
  \end{block}

\end{frame}

% ── Slide 8 — Stack Technique ────────────────────────────────────────────────
\begin{frame}{Stack Technique}

  \begin{columns}[T]

    \begin{column}{0.33\textwidth}
      \begin{block}{\faCode\quad Services métier}
        \begin{itemize}\small
          \item \textbf{Spring Boot} (Java)\\
                \scriptsize tous les microservices\normalsize
          \item \textbf{Spring Cloud Gateway}\\
                \scriptsize API Gateway\normalsize
          \item \textbf{Spring Cloud Eureka}\\
                \scriptsize service discovery\normalsize
          \item \textbf{Spring Cloud Config}\\
                \scriptsize configuration centralisée\normalsize
          \item \textbf{Keycloak}\\
                \scriptsize authentification JWT\normalsize
          \item \textbf{Apache Kafka}\\
                \scriptsize messaging asynchrone\normalsize
        \end{itemize}
      \end{block}
    \end{column}

    \begin{column}{0.33\textwidth}
      \begin{block}{\faDatabase\quad BDD \& Déploiement}
        \begin{itemize}\small
          \item \textbf{Base de données dédiée}\\
                \scriptsize par service\normalsize
          \item \textbf{Docker \& Docker Compose}\\
                \scriptsize conteneurisation\normalsize
          \item \textbf{Nginx}\\
                \scriptsize reverse proxy\normalsize
          \item \textbf{AWS EC2}\\
                \scriptsize déploiement cloud\normalsize
        \end{itemize}
      \end{block}
    \end{column}

    \begin{column}{0.33\textwidth}
      \begin{block}{\faChartBar\quad Observabilité}
        \begin{itemize}\small
          \item \textbf{Prometheus}\\
                \scriptsize métriques\normalsize
          \item \textbf{Loki}\\
                \scriptsize logs\normalsize
          \item \textbf{Tempo + OpenTelemetry}\\
                \scriptsize tracing distribué\normalsize
          \item \textbf{Grafana}\\
                \scriptsize visualisation unifiée\normalsize
        \end{itemize}
      \end{block}
    \end{column}

  \end{columns}

\end{frame}

% ── Slide 9a — Réalisation : Interface Utilisateur ───────────────────────────
\begin{frame}{Réalisation : Interface Utilisateur}

  \begin{columns}[T]
    \begin{column}{0.48\textwidth}
      \begin{center}
        \fbox{\parbox[c][0.45\textheight][c]{0.90\linewidth}{%
          \centering\color{primary!50}
          \faImage\\[6pt]
          \small [Screenshot Frontend 1]
        }}
      \end{center}
    \end{column}
    \begin{column}{0.48\textwidth}
      \begin{center}
        \fbox{\parbox[c][0.45\textheight][c]{0.90\linewidth}{%
          \centering\color{primary!50}
          \faImage\\[6pt]
          \small [Screenshot Frontend 2]
        }}
      \end{center}
    \end{column}
  \end{columns}

  \vspace{0.4em}
  \begin{block}{}
    \small Interface utilisateur permettant la \textbf{navigation}, la \textbf{gestion du panier}
    et le \textbf{passage de commandes}.
  \end{block}

\end{frame}

% ── Slide 9b — Réalisation : Backend & Services ───────────────────────────────
\begin{frame}{Réalisation : Backend \& Services}

  \begin{columns}[T]
    \begin{column}{0.48\textwidth}
      \begin{center}
        \fbox{\parbox[c][0.45\textheight][c]{0.90\linewidth}{%
          \centering\color{primary!50}
          \faImage\\[6pt]
          \small [Screenshot Backend 1]
        }}
      \end{center}
    \end{column}
    \begin{column}{0.48\textwidth}
      \begin{center}
        \fbox{\parbox[c][0.45\textheight][c]{0.90\linewidth}{%
          \centering\color{primary!50}
          \faImage\\[6pt]
          \small [Screenshot Backend 2]
        }}
      \end{center}
    \end{column}
  \end{columns}

  \vspace{0.4em}
  \begin{block}{}
    \small Vue sur les services enregistrés dans \textbf{Eureka} et les endpoints
    exposés via l'\textbf{API Gateway}.
  \end{block}

\end{frame}

% ── Slide 9c — Réalisation : Observabilité ───────────────────────────────────
\begin{frame}{Réalisation : Observabilité}

  \begin{center}
    \fbox{\parbox[c][0.38\textheight][c]{0.96\linewidth}{%
      \centering\color{primary!50}
      \faImage\\[6pt]
      [Screenshot Grafana Dashboard]
    }}
  \end{center}

  \vspace{0.3em}

  \begin{columns}[T]
    \begin{column}{0.33\textwidth}
      \begin{block}{\centering\small Prometheus}
        \centering\scriptsize Métriques collectées en temps réel
      \end{block}
    \end{column}
    \begin{column}{0.33\textwidth}
      \begin{block}{\centering\small Loki}
        \centering\scriptsize Logs centralisés et consultables
      \end{block}
    \end{column}
    \begin{column}{0.33\textwidth}
      \begin{block}{\centering\small Tempo}
        \centering\scriptsize Traces distribuées entre les services
      \end{block}
    \end{column}
  \end{columns}

\end{frame}

% ── Slide 10 — Défis & Solutions ─────────────────────────────────────────────
\begin{frame}{Défis \& Solutions}

  \begin{block}{\faTools\quad Complexité opérationnelle}
    \small \textbf{21 containers} à orchestrer simultanément. Les ressources nécessaires ont
    imposé un déploiement distribué sur \textbf{3 instances EC2} : bases de données, stack
    LGTM, Kafka et Zookeeper déportés du poste local.
  \end{block}

  \vspace{0.4em}

  \begin{block}{\faUserAlt\quad Charge de développement solo}
    \small \textbf{6 microservices} développés par une seule personne. Certaines
    fonctionnalités prévues n'ont pas pu être implémentées : orchestration
    \textbf{Kubernetes}, circuit breakers (\textbf{Resilience4j}).
  \end{block}

  \vspace{0.4em}

  \begin{block}{\faSyncAlt\quad Cohérence des transactions distribuées}
    \small Sans base de données partagée, les transactions multi-services posent un
    problème de cohérence. Solution : implémentation du \textbf{pattern SAGA} pour
    garantir la cohérence des données de manière fiable.
  \end{block}

\end{frame}

% ── Slide 13 — Ce que j'ai appris ────────────────────────────────────────────
\begin{frame}{Ce que j'ai appris}

  \begin{block}{\faBalanceScale\quad Le coût réel des microservices}
    \small L'\textbf{overhead opérationnel} des microservices est difficile à justifier :
    le seuil à partir duquel ils deviennent nécessaires est très élevé. Une PME n'en
    aura jamais besoin — il existe un cas d'usage très spécifique où ils surpassent
    le monolithe.
  \end{block}

  \vspace{0.35em}

  \begin{block}{\faPuzzlePiece\quad La modélisation des services est un art}
    \small Certains services ont été bien plus sollicités que d'autres, sans qu'il soit
    possible de le prévoir à l'avance. Avec le recul, certains microservices auraient
    pu être \textbf{fusionnés} en un seul service plus cohérent.
  \end{block}

  \vspace{0.35em}

  \begin{block}{\faEye\quad L'observabilité est non négociable}
    \small Une observabilité robuste est indispensable pour toute application, et dans
    le cas des microservices, le \textbf{tracing distribué est impératif}. Loki montre
    ses limites face à de grands volumes de logs — un remplacement par la stack
    \textbf{ELK} (Elasticsearch, Logstash, Kibana) aurait été plus adapté.
  \end{block}

\end{frame}

% ── Slide 14 — Conclusion & Perspectives ─────────────────────────────────────
\begin{frame}{Conclusion \& Perspectives}

  \begin{block}{\faTrophy\quad Bilan}
    \small \textbf{Alexandre} est une implémentation complète et fonctionnelle d'une
    architecture microservices moderne — service discovery, configuration centralisée,
    communication event-driven, observabilité complète, et déploiement cloud sur
    \textbf{AWS EC2}. Développé en solo dans un cadre académique, il démontre la
    faisabilité et la complexité réelle de ce type d'architecture.
  \end{block}

  \vspace{0.4em}

  \begin{columns}[T]
    \begin{column}{0.48\textwidth}
      \textbf{\color{primary}\small Fonctionnalités non implémentées}
      \begin{itemize}\small
        \item Orchestration \textbf{Kubernetes}
        \item Circuit breakers (\textbf{Resilience4j})
        \item Communication asynchrone généralisée
        \item Mise en cache (\textbf{Redis})
      \end{itemize}
    \end{column}
    \begin{column}{0.48\textwidth}
      \textbf{\color{primary}\small Améliorations possibles}
      \begin{itemize}\small
        \item Remplacement de Loki par la stack \textbf{ELK}
        \item \textbf{Autoscaling} horizontal des services
        \item Pipeline \textbf{CI/CD} automatisé
        \item Tests d'intégration inter-services
      \end{itemize}
    \end{column}
  \end{columns}

  \vspace{0.4em}

  \begin{alertblock}{}
    \centering\small
    Ce projet m'a permis de comprendre non seulement \textbf{comment} construire
    une architecture distribuée, mais surtout \textbf{quand} — et \textbf{quand
    ne pas} — la choisir.
  \end{alertblock}

\end{frame}

% ═══════════════════════════════════════════════════════════════════════════════
\end{document}
